% =============================================================================
% Capítulo 6 — Conclusão   (peso: leve)
% Fonte: plano §8. Contribuições, limitações (inclusive o que B3/B4 não respondem), trabalhos futuros.
% =============================================================================
\chapter{Conclusão}
\label{cap:conclusao}

\pendencia{Redigir por último.}

% -----------------------------------------------------------------------------
% ANOTAÇÃO PARA AS LIMITAÇÕES (não compila; retirada da Metodologia em
% 20/11/2026 por orientação do orientador). Ameaças à validade previstas, na
% taxonomia de Wohlin et al. (wohlin2024): rever com o que de fato ocorrer na
% coleta e converter em texto da seção de limitações. A Introdução (Seção
% Justificativa) também menciona o limite de generalização do Benchmark 2.
% -----------------------------------------------------------------------------
% As ameaças previstas seguem os quatro tipos de \citeonline{wohlin2024} e serão
% revistas ao fim da coleta, com o que de fato tiver ocorrido.
%
% \textbf{Validade interna.} O autor implementa as duas versões conhecendo a
% hipótese. Para reduzir esse viés no \textit{Benchmark}~1, as métricas primárias
% são estruturais e comportamentais, independentes de julgamento; o tempo, mais
% sujeito ao viés, é secundário, a ordem das trocas será declarada e a versão D
% passa por revisão independente antes da troca para verificar aderência
% idiomática ao SDK. O mesmo viés alcança a forma como a camada expõe as saídas
% estruturadas: ela será fixada em um ADR antes da implementação, e a alternativa
% preterida será discutida com os resultados. Como as trocas são feitas em
% sequência pela mesma pessoa, as últimas se beneficiam do aprendizado das
% primeiras e, na \versaoH, do adaptador já escrito; isso afeta o tempo, e não as
% métricas estruturais. A falta de equivalência funcional faria os dois
% experimentos compararem objetos diferentes, o que a suíte comum e a verificação
% simétrica antes de cada coleta evitam. Nas requisições agênticas, o descarte de
% execuções cujo número de chamadas de ferramenta diferir do esperado pode gerar
% viés de seleção pós-tratamento se diferir entre versões; por isso a taxa de
% descarte será registrada por versão. No \textit{Benchmark}~2, a leitura de
% energia pode ser contaminada por outros processos na GPU e variar entre sessões;
% as medidas são a exclusividade verificada a cada sessão, o congelamento de
% \textit{driver} e modelo, a execução em pares intercalados por sessão e o
% desconto da energia ociosa de cada sessão. Como a aplicação acessa o servidor
% pela rede, variações de rede e o desvio de sincronização de relógios entre as
% duas máquinas (potencialmente comparável ao tempo de processamento local da
% camada) afetam o alinhamento das janelas; o tempo local desconta a espera pelo
% servidor, o desvio de relógios é medido por sessão e o pareamento distribui a
% variação da rede entre as versões.
%
% \textbf{Validade externa.} O estudo usa uma única aplicação, de um único
% domínio, e deve ser lido como estudo de caso. O resultado do
% \textit{Benchmark}~1 depende do número de pontos de integração da aplicação, que
% será declarado. O cenário com recurso específico usa um único recurso (saídas
% estruturadas com esquema estrito) e não se estende, sem ressalva, a outros
% recursos proprietários (como ferramentas executadas no provedor ou
% \textit{streaming} heterogêneo). Há ainda assimetria de domínio entre os dois
% experimentos: o \textit{Benchmark}~1 avalia provedores comerciais em nuvem,
% enquanto o \textit{Benchmark}~2 avalia um único modelo aberto auto-hospedado em
% uma única GPU, o que não permite generalizar o efeito energético a outros
% modelos, outros equipamentos ou provedores comerciais. O \harness{} e os
% protocolos serão publicados para permitir a replicação.
%
% \textbf{Validade de construto.} As métricas estruturais medem o custo de
% substituição de forma indireta e podem não captar o esforço cognitivo de
% entender o código, limitação também reconhecida por \citeonline{mo2023}; por
% isso são combinadas com a métrica comportamental, com o tempo e com o registro
% separado do custo inicial de construção da camada (não incluído no limiar
% binário de troca). O consumo agregado de energia esconde a relação entre
% potência e duração, e por isso as duas serão reportadas separadamente (efeito de
% duração e efeito de potência). A telemetria NVML da GPU possui tolerância típica
% de sensor (cerca de $\pm 5\%$ em GPUs de alta potência) e leituras de potência
% podem refletir janelas móveis suavizadas, razão pela qual se adota o contador
% acumulado de energia; além disso, ela não alcança o trabalho que a camada
% executa na CPU da máquina da aplicação, parcela estimada por um limite superior
% que tende a superestimar o consumo real. O efeito energético é medido apenas no
% cenário comum; o confronto do cenário com saídas estruturadas reutiliza essa
% mesma medida como aproximação.
%
% \textbf{Validade de conclusão.} Com amostras pequenas em janelas curtas de GPU
% ou alta variação residual de comprimento de saída, o poder estatístico pode ser
% insuficiente para demonstrar não inferioridade dentro da margem $\Delta_0 = 5\%$,
% motivando o piloto e a contingência energética da atividade 10. Além disso, a
% margem de 5\% é uma decisão pré-registrada ancorada na literatura e na
% resolução da telemetria, mas outras margens alterariam o veredito binário; por
% isso os intervalos de confiança, os tamanhos de efeito (delta de Cliff e
% deslocamento relativo mediano) e a correção de Holm-Bonferroni serão reportados
% integralmente ao lado dos valores-$p$.
